\documentclass[11pt,a4paper]{article}
\usepackage[utf8]{inputenc}
\usepackage[T1]{fontenc}
\usepackage{CJKutf8}
\usepackage{geometry}
\geometry{margin=2.2cm}
\usepackage{booktabs}
\usepackage{array}
\usepackage{longtable}
\usepackage{hyperref}
\usepackage{xcolor}
\usepackage{listings}
\usepackage{siunitx}
\usepackage{caption}

\hypersetup{
  colorlinks=true,
  linkcolor=blue!50!black,
  urlcolor=blue!60!black,
  citecolor=blue!50!black
}

\lstset{
  basicstyle=\ttfamily\footnotesize,
  breaklines=true,
  columns=fullflexible,
  frame=single,
  framerule=0.3pt,
  showstringspaces=false
}

\newcommand{\NEB}{NEBULA}
\newcommand{\NPL}{NEBULA-Plus}

\title{\NEB{} and \NPL{} geometry, performance, and data-source notes\\
       \large (collected for the kondo\_smsimulator / smsimulator5.5 project)}
\author{Notes compiled by Baiting Tian}
\date{2026-05-15}

\begin{document}
\begin{CJK}{UTF8}{gbsn}
\maketitle

\section{Layout and orientation (executive summary)}

Along the beam direction, from target downstream:
\[
\text{target}  \;\longrightarrow\;
\textbf{\NPL{} (front)} \;\longrightarrow\; \textbf{\NEB{} (back)}.
\]

That is, \textbf{\NPL{} sits in front of \NEB{}} (between FDC2/HOD and \NEB{}).
\NPL{} contains \textbf{90 plastic-scintillator neutron bars}, i.e.\
\textbf{30 bars fewer} than \NEB{}'s 120 bars. The plastic bar dimensions
($\SI{120}{mm}\times\SI{1800}{mm}\times\SI{120}{mm}$) and the veto layout (\SI{10}{mm} plastic
scintillator at the upstream face of each wall) are otherwise identical to \NEB{}.

\NPL{} was built by a Japan / LPC-Caen collaboration under the French ANR EXPAND project,
shipped to the RIBF in summer 2022 and mechanically installed in October 2022.
Initial in-beam commissioning with $^{7}\mathrm{Li}(p,n)^{7}\mathrm{Be}$ was scheduled
for the next \NEB{} campaign~\cite{aprApr2022}.

\NPL{} is now \textbf{implemented in smsimulator5.5} (branch
\texttt{nebula-plus-impl} merged to \texttt{main} on 2026-05-16, tags
\texttt{nebula-plus-phase1/2/3}). End-to-end verification ran 100
\SI{200}{MeV} pencil-beam neutrons through both detectors and recovered
neutron candidates from three reco classes (\NEB{}-only, \NPL{}-only,
joint); the joint reco produced $\sim\!\!\times 1.86$ as many candidates
as \NEB{} alone, consistent with the qualitative ``$\sim\!\!\times 2$''
gain quoted in Tanaka 2024~\cite{tanaka2024}. See \S\ref{sec:impl} and
\S\ref{sec:provenance} for full implementation details and the source of
every authoritative parameter.

\section{Internal geometry}

\subsection{\NEB{} (existing, already in smsimulator5.5)}

Two walls (Wall-A, Wall-B), each with two sub-layers of neutron bars
and one upstream veto layer:

\begin{table}[h]
\centering
\caption{\NEB{} layout (s021 configuration; coordinates from the smsimulator5.5
\texttt{geometry/s021/NEBULA\_Detectors\_samurai21.csv}).}
\begin{tabular}{lccc}
\toprule
Element & \# bars & $Z$ rel.\ \NEB{} origin (mm) & $X$ range (mm) \\
\midrule
Wall-A Veto    & 12 & $-275,\;-260$         & $-$ \\
Wall-A Sub1    & 30 & $0$                   & $+1908 \to -1641.6$\\
Wall-A Sub2    & 30 & $+130$                & same\\
Wall-B Veto    & 12 & $+575.3,\;+590.3$     & $-$\\
Wall-B Sub1    & 30 & $+850.3$              & same\\
Wall-B Sub2    & 30 & $+980.3$              & same\\
\midrule
\NEB{} origin in target frame (s021) & \multicolumn{3}{l}{$Z = \SI{9784}{mm}$}\\
Bar size (Neut)   & \multicolumn{3}{l}{$120 \times 1800 \times \SI{120}{mm}$ ($W\times H\times L$)}\\
Bar pitch in $X$  & \multicolumn{3}{l}{$\SI{122.4}{mm}$ (i.e.\ \SI{2.4}{mm} air gap)}\\
Bar centroid $Y$  & \multicolumn{3}{l}{$\SI{-35.4}{mm}$ (slightly below beam axis;
                       see \texttt{NEBULA\_Pos\_FromMagCenter.csv})}\\
Veto size         & \multicolumn{3}{l}{$320 \times 1900 \times \SI{10}{mm}$}\\
Time resolution per PMT & \multicolumn{3}{l}{$\sigma_t = \SI{0.240}{ns}$
                       (\texttt{NEBULA\_samurai21.csv})}\\
Bar material (in current sim) & \multicolumn{3}{l}{\texttt{G4\_PLASTIC\_SC\_VINYLTOLUENE}
                       (NIST PVT analog of BC408/EJ200)}\\
\bottomrule
\end{tabular}
\end{table}

Total: $120$ Neut bars $+\,24$ Veto bars $=\,144$ counters.

\subsection{\NPL{} (now implemented in smsimulator5.5)}
\label{sec:internal-npl}

From the LPC-Caen / RIKEN nptool plugin
\texttt{commissioning/detector.yaml} (see \S\ref{sec:sources}), the four
sub-layers of \NPL{} are:

\begin{table}[h]
\centering
\caption{\NPL{} layout from the nptool plugin's \texttt{detector.yaml}.
$Z$ is given in the \NPL{}-internal frame (its own Wall-A Sub1 is the origin).
$X_{0}$ is the centre of the leftmost bar; the rest of the row follows
$X_i = X_0 + i \times \SI{120}{mm}$.}
\begin{tabular}{lccccc}
\toprule
Wall & SubLayer & \# bars & $X_0$ (mm) & $X$ range (mm) & $Z$ (mm)\\
\midrule
A & Sub1 & 22 & $-1320$ & $-1320 \to +1200$ & $0$\\
A & Sub2 & 22 & $-1320$ & $-1320 \to +1200$ & $130$\\
B & Sub1 & 23 & $-1320$ & $-1320 \to +1320$ & $845$\\
B & Sub2 & 23 & $-1320$ & $-1320 \to +1320$ & $975$\\
\midrule
\multicolumn{6}{l}{Bar size (Neut): $120 \times 1800 \times \SI{120}{mm}$ — identical to \NEB{}}\\
\multicolumn{6}{l}{Bar material: BC400 in nptool (essentially equivalent to \NEB{}'s PVT analog)}\\
\multicolumn{6}{l}{Bar pitch in $X$: \SI{120}{mm} (no air gap in the plugin model;
                  \texttt{InterModule = 1\,mm} declared but unused)}\\
\multicolumn{6}{l}{Veto size: $320 \times 1900 \times \SI{10}{mm}$}\\
\multicolumn{6}{l}{Veto distance to wall: \SI{100}{mm} upstream of the corresponding Neut wall
                  (\texttt{WallToVeto = 10\,cm})}\\
\multicolumn{6}{l}{Veto count: $12$/wall (the plugin selects $12$ when $N_\mathrm{bars}>15$,
                  else $6$)}\\
\multicolumn{6}{l}{PMT: Hamamatsu H11284 assemblies; FASTER digital electronics
                  developed at LPC-Caen}\\
\bottomrule
\end{tabular}
\end{table}

Total: $22+22+23+23 = 90$ Neut bars + $24$ Veto bars = $\mathbf{114}$ counters,
i.e.\ \textbf{30 Neut bars fewer than \NEB{}}.

\subsection{Absolute $Z$ of the \NPL{} origin}

The \texttt{detector.yaml} files in the LPC-Caen nptool plugin use a placeholder
target block ($x=1,\;y=2,\;z=\SI{3}{mm}$); the absolute $Z$ of the \NPL{} origin
relative to the SAMURAI target is therefore \emph{not} captured in any of the
config files we have read-access to. The home directories of the responsible
authors (\texttt{/home/kondo}, \texttt{/home/otsu}, \texttt{/home/gibelin},
\texttt{/home/jlemarie}) on \texttt{ribfana04} were permission-denied.

\textbf{Chosen value used in smsimulator5.5 (placeholder; survey-grade value still
to be obtained from LPC-Caen):}
\[
Z_{\NPL{},\,\text{origin}} = \SI{8089}{mm}\quad\text{(target frame)},
\]
selected so the Wall-A--Sub1 spacings $845 \to 850 \to 850.3$ mm walk uniformly
from \NPL{} into \NEB{}:

\begin{table}[h]
\centering
\caption{\NPL{}+\NEB{} target-frame layout used in smsimulator5.5
(combined of \S\ref{sec:internal-npl} numbers and chosen $Z_0$).}
\begin{tabular}{lrr}
\toprule
Wall / sub-layer & $Z$ (mm) & Sub1$\to$Sub1 spacing (mm)\\
\midrule
NPL Wall-A Veto                 & 7989  & --\\
\textbf{NPL Wall-A Sub1}        & \textbf{8089}  & --\\
NPL Wall-A Sub2                 & 8219  & --\\
NPL Wall-B Veto                 & 8834  & --\\
\textbf{NPL Wall-B Sub1}        & \textbf{8934}  & $845.0$ (NPL internal)\\
NPL Wall-B Sub2                 & 9064  & --\\
NEB Wall-A Veto                 & 9509  & --\\
\textbf{NEB Wall-A Sub1}        & \textbf{9784}  & $850.0$ (gap chosen)\\
NEB Wall-A Sub2                 & 9914  & --\\
NEB Wall-B Veto                 & $10359.3$ & --\\
\textbf{NEB Wall-B Sub1}        & $\mathbf{10634.3}$ & $850.3$ (NEB internal)\\
NEB Wall-B Sub2                 & $10764.3$ & --\\
\bottomrule
\end{tabular}
\end{table}

The chosen \SI{850}{mm} NPL-B$\to$NEB-A gap can be overridden at runtime via
\texttt{/samurai/geometry/NEBULAPlus/Position 0\;0\;<Z>\;mm} (and similarly
for NEB) once a real survey is provided.

\section{Material / electronics summary}

\begin{itemize}
  \item Plastic: BC400 (\NPL{}) and BC408/EJ200 (\NEB{}) are PVT-based
        scintillators; densities and light yields agree to better than
        $1\,\%$. In the Geant4 simulation both can be represented by
        \texttt{G4\_PLASTIC\_SC\_VINYLTOLUENE}.
  \item Refractive index $n = 1.58$; bulk attenuation length \SI{6680}{mm}
        (declared in \texttt{NebulaPlusGeant4.cxx}).
  \item PMTs: \NPL{} uses Hamamatsu H11284; \NEB{} uses Hamamatsu H7195 or
        equivalent (depending on bar batch).
  \item DAQ: \NPL{} uses the FASTER digital readout developed at LPC-Caen,
        integrated with SAMURAI via a coincidence trigger derived from the
        two-PMT signal pair of each bar~\cite{aprApr2022}.
\end{itemize}

\section{Detection efficiency}

\subsection{\NEB{} alone}

Single-neutron detection efficiency at $E_n = \SI{200}{MeV}$:
\[
\varepsilon_{1n}^{\NEB{}} = 32.5 \pm 0.3\,(\text{stat}) \pm 0.9\,(\text{syst})\,\%,
\]
measured with the $^{7}\mathrm{Li}(p,n)^{7}\mathrm{Be}$ reaction
(Kondo et al., NIM B 463 (2020) 173--178~\cite{kondo2020nim}).

\subsection{\NEB{} $\cup$ \NPL{} (combined)}

The \NEB{}+\NPL{} array effectively doubles the number of layers
(four neutron sub-layers $\to$ eight). The Tanaka et al.\ review
(arXiv:2412.17887~\cite{tanaka2024}) states, qualitatively:

\begin{itemize}
  \item \textbf{1-neutron} efficiency: increases by \textbf{nearly a factor of 2}
        (i.e.\ from $\sim 32.5\,\%$ to $\sim 60\text{--}65\,\%$ at \SI{200}{MeV},
        rough extrapolation).
  \item \textbf{4-neutron} efficiency: improves by \textbf{almost one order of
        magnitude} relative to \NEB{} alone.
\end{itemize}

Absolute combined numbers and the energy dependence have to come from the
in-beam commissioning data (\NPL{} alone $\varepsilon_{1n}$ to be determined
once the $^{7}\mathrm{Li}(p,n)$ run is analysed --- planned but not yet
published as of the APR~2022 report).

\section{Local file paths (where this information lives)}
\label{sec:sources}

\subsection{On this workstation}

\begin{description}
  \item[smsimulator5.5 NEBULA G4 module]
    \texttt{kondo\_smsimulator/sources/smg4lib/detectors/NEBULA/}
  \item[NEBULA geometry CSVs (s021)]
    \texttt{kondo\_smsimulator/work/geometry/s021/NEBULA\_samurai21.csv} \\
    \texttt{kondo\_smsimulator/work/geometry/s021/NEBULA\_Detectors\_samurai21.csv}
  \item[anaroot NEBULA parameter db]
    \texttt{anaroot/db/NEBULA.20250625.xml} (and dated siblings) \\
    \texttt{anaroot/db/NEBULA\_Pos\_FromMagCenter.csv} \\
    \texttt{anaroot/db/NEBULA.ods}
  \item[anaroot \NPL{} placeholder (only veto names V301--V410 at $Z=0$)]
    \texttt{anaroot/db/NEBULA.20250625.xml}\\
    \texttt{anaroot/src/sources/Reconstruction/SAMURAI/src/TArtSAMURAIParameters.cc:544}
    \\ (comment: \texttt{// for new config. with NEBULA Plus})
  \item[Papers (downloaded)]
    \texttt{docs/papers/NEBULA-Plus\_Kondo\_RIKEN\_APR2022.pdf} \\
    \texttt{docs/papers/NEBULA-Plus\_review\_Tanaka2024\_arxiv2412.17887.pdf} \\
    \texttt{docs/papers/multi-neutron\_SAMURAI\_Yang2019\_arxiv1903.11256.pdf} \\
    \texttt{docs/papers/NEBULA\_RIBF\_factsheet.pdf} \\
    \texttt{/docs/papers/NEBULA\_indico\_talk.pdf}
\end{description}

\subsection{On the analysis server \texttt{ribfana04}}

These are the definitive \NPL{} geometry and code locations on the s053
experiment account (read-access only):

\begin{description}
  \item[nptool plugin (Geant4 + Physics)]
    \texttt{/home/s053/exp/exp2301\_s053/nptool/plugin/nebula-plus/sources/src/} \\
    Main files:
    \begin{itemize}
      \item \texttt{NebulaPlusGeant4.cxx}, \texttt{NebulaPlusGeant4.h}
            --- Geant4 module / veto construction, materials, vis attributes.
      \item \texttt{NebulaPlusDetector.cxx}, \texttt{NebulaPlusDetector.h}
            --- yaml parsing, AddLayer / AddSBT / AddSBV; Z is relative to target.
      \item \texttt{NebulaPlusPhysics.h}, \texttt{NebulaPlusData.h/cxx},
            \texttt{NebulaPlusSpectra.cxx}, \texttt{NebulaPluslinkdef.hh}.
    \end{itemize}
  \item[detector configurations]
    \texttt{/home/s053/exp/exp2301\_s053/nptool/commissioning/detector.yaml}
    (the only one with \texttt{X\_start = -1320\,mm}; \texttt{s053/detector.yaml}
    and \texttt{sandbox/detector.yaml} use \texttt{X = 0}.)
  \item[Calibrations]
    \texttt{/home/s053/exp/exp2301\_s053/nptool/commissioning/calibration/}\\
    \quad \texttt{NebulaPlus\_Energy.cal},
    \texttt{NebulaPlus\_Up\_Energy.cal},
    \texttt{NebulaPlus\_Down\_Energy.cal},\\
    \quad \texttt{NebulaPlus\_Tdiff\_Offset.cal},
    \texttt{NebulaPlus\_Tdiff\_Offset\_Jan2023.cal}
  \item[Online monitor (s053)]
    \texttt{/home/s053/exp/exp2301\_s053/sm\_online\_utils/s053/chknebulaplus/OnlineMonitor.cc}
    (counts/rates only; no geometry constants).
  \item[Raw data + FASTER setup files]
    \texttt{/home/s053/rawdata/nebulaplus/NEBULA\_Test\_2025*.fast} (+\texttt{.setup})
  \item[FASTER electronics source tree]
    \texttt{/home/s053/exp/exp2301\_s053/faster/faster/fasterac-2.18/}
  \item[Rsync from \NPL{} DAQ host]
    \texttt{/home/s053/exp/exp2301\_s053/scripts/rsync\_expand\_test.sh}
    \\ pulls from \texttt{expand@expandacq:./Projects/test\_veto/data/Test\_VETO/} into
    \texttt{\$HOME/rawdata/nebulaplus/}.
\end{description}

\subsection{External}

\begin{itemize}
  \item Project page (LPC-Caen):
    \url{https://nebula-plus.in2p3.fr/}
    (with \texttt{/manual/} and \texttt{/documentation-and-private-files/}
    sub-areas behind login).
  \item RIKEN APR 56 (2023) p.\,91 -- installation report:
    \url{https://www.nishina.riken.jp/researcher/APR/APR056/pdf/91.pdf}
\end{itemize}

\section{Implementation in smsimulator5.5}
\label{sec:impl}

Implemented on branch \texttt{nebula-plus-impl} (merged to \texttt{main}
2026-05-16, tags \texttt{nebula-plus-phase1/2/3} mark each phase boundary).
The corresponding spec and plan are
\texttt{docs/superpowers/specs/2026-05-15-nebula-plus-design.md}
and
\texttt{docs/superpowers/plans/2026-05-15-nebula-plus.md}.

\subsection{Build environment}

\begin{lstlisting}[language=bash]
# 1. Activate the anaroot conda environment (Geant4 11.x, ROOT 6.36,
#    Xerces-C, anaroot -- all linked to that env)
micromamba activate anaroot-env

# 2. spdlog must be installed in the env (used by smlogger). One-time:
micromamba install -n anaroot-env -c conda-forge spdlog

# 3. Source project setup (defines SMSIMDIR, TARTSYS, paths)
cd $HOME/workspace/dpol/smsimulator5.5
source setup.sh

# 4. Configure + build
mkdir -p build && cd build
cmake .. -DCMAKE_BUILD_TYPE=Release \
         -DBUILD_TESTS=ON -DBUILD_APPS=ON \
         -DBUILD_ANALYSIS=ON -DWITH_ANAROOT=ON
cmake --build . -j$(nproc)
\end{lstlisting}

Important: \texttt{TARTSYS} must be set \emph{before} \texttt{cmake ..} runs;
otherwise \texttt{WITH\_ANAROOT} silently flips off and the
\texttt{NEBULASimDataConverter\_TArtNEBULAPla} translation unit fails to find
\texttt{TArtNEBULAPla.hh}. \texttt{setup.sh} on host \texttt{TBTdebian} exports
\texttt{TARTSYS=/home/tian/software/anaroot/install}.

\subsection{Geometry configuration (CSV)}

Two pairs of CSV files drive the geometry. The smsim G4 messengers read them
at runtime via \texttt{/samurai/geometry/.../ParameterFileName}.

\begin{table}[h]
\centering
\caption{Geometry-CSV inventory under
\texttt{smsimulator5.5/configs/simulation/geometry/}. Each pair is
\{global-parameters file, per-bar table\}.}
\begin{tabular}{lll}
\toprule
Detector & Global params & Per-bar table\\
\midrule
\NEB{} (Dayone)    & \texttt{NEBULA\_Dayone.csv}           & \texttt{NEBULA\_Detectors\_Dayone.csv}\\
\NPL{} (s021)      & \texttt{s021/NEBULAPlus\_samurai21.csv}
                                                          & \texttt{s021/NEBULAPlus\_Detectors\_samurai21.csv}\\
\bottomrule
\end{tabular}
\end{table}

The \NPL{} CSV pair was generated from the numbers in
\S\ref{sec:internal-npl}. The global-params file content is
shown below verbatim; the per-bar file has 114 rows ($90+24$) following the
\texttt{Layer/SubLayer/PositionX/Y/Z} schema also used by \NEB{}.

\begin{lstlisting}
% configs/simulation/geometry/s021/NEBULAPlus_samurai21.csv
Position, 0, 0, 8089,
NeutSize, 120, 1800, 120,
VetoSize, 320, 1900, 10,
Angle, 0, 0, 0,
////input time reso for one PMT
TimeReso, 0.240,
\end{lstlisting}

\subsection{Geant4 module classes (\texttt{libs/smg4lib/})}

\NPL{}'s G4 module mirrors \NEB{}'s file-for-file. The construction-side
classes live in \texttt{src/construction/\{src,include\}/}; the data-side
classes live in \texttt{src/data/\{src,include\}/} with mirror public headers
also in \texttt{libs/smg4lib/include/}.

\begin{table}[h]
\centering
\caption{G4-side class inventory. \NPL{} files were created by sed-clone
of the corresponding \NEB{} files in branch \texttt{nebula-plus-impl}
Phase 1.}
\begin{tabular}{lll}
\toprule
Role & \NEB{} class & \NPL{} class\\
\midrule
G4 geometry         & \texttt{NEBULAConstruction}            & \texttt{NEBULAPlusConstruction}\\
G4 UI messenger     & \texttt{NEBULAConstructionMessenger}   & \texttt{NEBULAPlusConstructionMessenger}\\
Sensitive detector  & \texttt{NEBULASD}                      & \texttt{NEBULAPlusSD}\\
CSV reader          & \texttt{NEBULASimParameterReader}      & \texttt{NEBULAPlusSimParameterReader}\\
SimData branch reg. & \texttt{NEBULASimDataInitializer}      & \texttt{NEBULAPlusSimDataInitializer}\\
Step$\to$hit conv.  & \texttt{NEBULASimDataConverter\_TArtNEBULAPla}
                                                             & \texttt{NEBULAPlusSimDataConverter\_TArtNEBULAPlusPla}\\
Run-param ROOT cls  & \texttt{TNEBULASimParameter}           & \texttt{TNEBULAPlusSimParameter}\\
Hit ROOT cls        & \texttt{TArtNEBULAPla} (anaroot)       & \texttt{TArtNEBULAPlusPla} (smsim, \textbf{new})\\
\bottomrule
\end{tabular}
\end{table}

Both detectors are instantiated and attached to the SAMURAI experimental
hall in \texttt{libs/sim\_deuteron\_core/src/DeutDetectorConstruction.cc}
(NPL added at Phase 1 task P1.11). The corresponding SimData initializer
and converter are registered in \texttt{apps/sim\_deuteron/main.cc}
(P1.12).

\subsection{Runtime \texttt{messenger} commands}

\begin{lstlisting}
# Load both detectors (works in any .mac executed by bin/sim_deuteron)
/samurai/geometry/NEBULA/ParameterFileName           {SMSIMDIR}/configs/simulation/geometry/NEBULA_Dayone.csv
/samurai/geometry/NEBULA/DetectorParameterFileName   {SMSIMDIR}/configs/simulation/geometry/NEBULA_Detectors_Dayone.csv

/samurai/geometry/NEBULAPlus/ParameterFileName       {SMSIMDIR}/configs/simulation/geometry/s021/NEBULAPlus_samurai21.csv
/samurai/geometry/NEBULAPlus/DetectorParameterFileName  {SMSIMDIR}/configs/simulation/geometry/s021/NEBULAPlus_Detectors_samurai21.csv

# Optionally override the global Z when a real LPC survey lands
/samurai/geometry/NEBULAPlus/Position 0 0 8089 mm

/samurai/geometry/Update
\end{lstlisting}

A minimal smoke macro is shipped at
\texttt{configs/simulation/macros/test\_nebula\_plus.mac}.

\subsection{Output ROOT-tree branches}

Each event of \texttt{bin/sim\_deuteron} produces (relevant branches only):

\begin{description}
  \item[\texttt{NEBULAPla}]     \texttt{TClonesArray} of
        \texttt{TArtNEBULAPla} (anaroot class) --- \NEB{} hits, with
        \texttt{Layer} = 3 (wall A) or 4 (wall B).
  \item[\texttt{NEBULAPlusPla}] \texttt{TClonesArray} of
        \texttt{TArtNEBULAPlusPla} (smsim class) --- \NPL{} hits, with
        \texttt{Layer} = 1 (wall A) or 2 (wall B).
  \item[\texttt{TNEBULASimParameter}, \texttt{TNEBULAPlusSimParameter}]
        per-run snapshots of the detector parameter map.
\end{description}

\subsection{Reconstruction (\texttt{libs/analysis/})}

The legacy \texttt{NEBULAReconstructor} was refactored in Phase 2 into a
small hierarchy:

\begin{description}
  \item[\texttt{NEBULABaseReco}] abstract base. Holds public algorithm
        steps shared by all derived classes:
        \texttt{ClusterHits},
        \texttt{ReconstructFromCluster},
        \texttt{CalculateBeta},
        \texttt{CalculateNeutronEnergy},
        \texttt{ApplyPositionSmearing},
        \texttt{ApplyTimeSmearing}.
        Pure virtual: \texttt{ExtractHits()}.
  \item[\texttt{NEBULAReco}] single-detector \NEB{}; \texttt{ExtractHits}
        reads \texttt{TClonesArray<TArtNEBULAPla>}. Preserves the legacy
        behaviour bit-for-bit (regression-pinned by
        \texttt{test\_NEBULAReco\_legacy\_compat}).
  \item[\texttt{NEBULAPlusReco}] single-detector \NPL{}; \texttt{ExtractHits}
        reads \texttt{TClonesArray<TArtNEBULAPlusPla>}, tags
        \texttt{wall\_tag} 1/2.
  \item[\texttt{NebulaJointReco}] takes both detector inputs and merges
        them through a single \texttt{ExtractHits} call; cluster\-er sees
        all 8 sub-layers and tags \texttt{wall\_tag} 1--4.
\end{description}

Default parameter values (preserved from the legacy class):
\texttt{fTimeWindow=10.0\,ns}, \texttt{fEnergyThreshold=1.0\,MeV},
\texttt{fPositionSmearing=5.0\,mm}, \texttt{fTimeSmearing=0.5\,ns}.

The reco \emph{algorithms} are now independent of anaroot's
\texttt{TArtRecoNeutron} / \texttt{TArtCalibNEBULA}. Anaroot is still a
build-time dependency because the \NEB{}-hit data class
\texttt{TArtNEBULAPla} lives in anaroot and is used as the on-disk schema
for the \texttt{NEBULAPla} branch. Removing that dependency is left as a
follow-up (would require introducing a smsim-side mirror class for
\NEB{}, plus a schema-evolution step for existing \texttt{.root} files).

\subsection{End-to-end verification (2026-05-16)}

Ran 100 monoenergetic \SI{200}{MeV} neutrons through the
\NPL{}$+$\NEB{} geometry; see \texttt{tests/integration/test\_sim\_deuteron\_with\_nebula\_plus.sh}.

\begin{table}[h]
\centering
\caption{End-to-end sim$+$reco numbers (100 incident \SI{200}{MeV}
neutrons, pencil beam from $Z=-2000$ mm).}
\begin{tabular}{lrr}
\toprule
Stage / class & Events with hits/candidates & Total entries\\
\midrule
\NEB{}      hits in tree     & 51 & 159 hits\\
\NPL{}      hits in tree     & 55 & 227 hits\\
\midrule
\texttt{NEBULAReco}        candidates & 32 & 35\\
\texttt{NEBULAPlusReco}    candidates & 39 & 42\\
\texttt{NebulaJointReco}   candidates & 55 & 65\\
\bottomrule
\end{tabular}
\end{table}

The joint-to-NEBULA candidate ratio is $65/35\approx 1.86$, consistent with
the qualitative ``$\sim\!\times 2$'' single-neutron gain quoted in the
Tanaka 2024 review~\cite{tanaka2024}.

\subsection{Tests}

\begin{table}[h]
\centering
\caption{New tests added by Phases 1--3 (\texttt{ctest} labels in
parentheses).}
\begin{tabular}{ll}
\toprule
File & Cases\\
\midrule
\texttt{tests/unit/test\_NEBULAPlusSimParameterReader.cc}    (unit)         & 3 \\
\texttt{tests/unit/test\_NEBULAPlusReco\_basic.cc}            (unit)         & 3 \\
\texttt{tests/unit/test\_NebulaJointReco\_combined.cc}        (unit)         & 3 \\
\texttt{tests/analysis/test\_NEBULAReco\_legacy\_compat.cc}   (analysis)     & 1 \\
\texttt{tests/integration/test\_sim\_deuteron\_with\_nebula\_plus.sh} (integration) & 1 \\
\bottomrule
\end{tabular}
\end{table}

After Phases 1--3, the project ctest count went from
\textbf{18 passing / 20 failing} to \textbf{30 passing / 20 failing}
($+12$ passing, $0$ regression on previously-passing tests). The 20
pre-existing failures (\texttt{ParticleTrajectoryTest.*},
\texttt{TargetReconstructorTest.*}) are unrelated to \NPL{} and remain
out-of-scope here.

\section{Provenance: where each NPL number came from}
\label{sec:provenance}

\begin{table}[h]
\centering
\caption{Information sources for every authoritative \NPL{} parameter used
in smsimulator5.5.}
\begin{tabular}{p{4.8cm}lp{6.5cm}}
\toprule
Parameter & Value & Source\\
\midrule
Bar dimensions ($W\times H\times L$)
  & $120\times 1800\times \SI{120}{mm}$
  & \texttt{NebulaPlusGeant4.cxx} (LPC plugin)\\
Veto dimensions
  & $320\times 1900\times \SI{10}{mm}$
  & \texttt{NebulaPlusGeant4.cxx} (LPC plugin)\\
Number of Neut bars per sub-layer
  & 22, 22, 23, 23
  & \texttt{commissioning/detector.yaml} (LPC plugin)\\
Number of Veto bars per wall
  & 12
  & \texttt{NebulaPlusGeant4.cxx} (\texttt{m\_NbrModule>15}$\Rightarrow$12 branch)\\
Bar pitch in $X$
  & \SI{120}{mm} (touching)
  & \texttt{NebulaPlusDetector::AddLayer} steps by \texttt{m\_BarWidth = 120}\\
Veto pitch in $X$
  & \SI{310}{mm}
  & NEBULA convention reused (same Veto type)\\
Veto distance to wall
  & \SI{100}{mm} upstream of Neut Sub1
  & \texttt{NebulaPlusGeant4.cxx::WallToVeto}\\
Sub-layer Z gap
  & \SI{130}{mm}
  & \texttt{commissioning/detector.yaml} (entries at Z=0/130/845/975)\\
$X$-start of first bar
  & $-\SI{1320}{mm}$ (asymmetric)
  & \texttt{commissioning/detector.yaml} (\texttt{pos: -132 0 0 cm})\\
Bar material (G4)
  & \texttt{G4\_PLASTIC\_SC\_VINYLTOLUENE}
  & Same as \NEB{} in smsimulator5.5; sufficient PVT analog of BC400/EJ200\\
Refractive index $n$
  & 1.58
  & \texttt{NebulaPlusGeant4.cxx::MaterialIndex}\\
Bulk attenuation length
  & \SI{6680}{mm}
  & \texttt{NebulaPlusGeant4.cxx::Attenuation}\\
PMT time resolution per side
  & \SI{0.240}{ns}
  & NEBULA value reused (no LPC commissioning number yet)\\
PMT model
  & Hamamatsu H11284
  & APR 56 (2023) p.\,91~\cite{aprApr2022}\\
Electronics
  & FASTER (LPC-Caen)
  & APR 56 (2023) p.\,91~\cite{aprApr2022}\\
\midrule
\textbf{Absolute $Z$ of NPL origin (target frame)}
  & \textbf{\SI{8089}{mm}}
  & \textbf{Chosen for smsim5.5} (uniform Sub1$\to$Sub1 spacing 845/850/850.3
        mm across NPL-A$\to$NPL-B$\to$NEB-A$\to$NEB-B; survey-grade value still
        to be obtained from LPC-Caen)\\
Layer index assignment
  & NPL: L=1,2; NEB: L=3,4
  & anaroot db convention (\texttt{V301-V310, V401-V410} pre-reserved;
        comment in \texttt{TArtSAMURAIParameters.cc:544})\\
Bar/Veto ID ranges
  & 101--122, 201--222, 301--323, 401--423; 501--512, 601--612
  & Chosen for smsim5.5 (Layer$\times$100 + bar index;
        Veto distinguished by 500/600 prefix and \texttt{SubLayer=0})\\
Single-neutron efficiency gain
  & $\sim\!\times 2$ vs \NEB{} alone
  & Tanaka 2024~\cite{tanaka2024} (arXiv:2412.17887; PTEP, to appear).
        Confirmed at $\times 1.86$ by the smsim end-to-end run above.\\
Four-neutron efficiency gain
  & $\sim\!\times 10$ vs \NEB{} alone
  & Tanaka 2024~\cite{tanaka2024}; not yet validated in smsim
        (requires a multi-neutron breakup macro)\\
NPL installation date
  & October 2022
  & APR 56 (2023) p.\,91~\cite{aprApr2022}\\
\bottomrule
\end{tabular}
\end{table}

For NEBULA-specific numbers (Z=9784 mm origin, 122.4 mm $X$ pitch, 35.4 mm
$Y$ offset, 0.240 ns time resolution), the authoritative source is the
existing CSV \texttt{configs/simulation/geometry/NEBULA\_Dayone.csv} and
the anaroot db files listed in \S\ref{sec:sources}; the design paper is
Kondo 2020 NIM B~\cite{kondo2020nim}.

\begin{thebibliography}{9}
\bibitem{kondo2020nim} Y.~Kondo, T.~Tomai, T.~Nakamura,
``Recent progress and developments for experimental studies with the SAMURAI
spectrometer,'' Nucl.\ Instrum.\ Methods Phys.\ Res.\ B \textbf{463} (2020) 173--178.
DOI: \texttt{10.1016/j.nimb.2019.04.085}.
\bibitem{aprApr2022} N.~A.~Orr, J.~Gibelin, Y.~Kondo, H.~Otsu, on behalf of the
\NPL{} collaboration, ``Installation of the \NPL{} neutron array,'' RIKEN Accel.\
Prog.\ Rep.\ \textbf{56} (2023) p.\,91.
\bibitem{tanaka2024} J.~Tanaka et al., ``Detectors for next-generation
quasi-free scattering experiments at the RIBF,'' arXiv:2412.17887 (to appear in PTEP).
\bibitem{nakamura2016} T.~Nakamura, Y.~Kondo, ``Large acceptance spectrometers for
invariant mass spectroscopy of exotic nuclei and future developments,''
Nucl.\ Instrum.\ Methods Phys.\ Res.\ B \textbf{376} (2016) 156. DOI:
\texttt{10.1016/j.nimb.2016.01.003}.
\end{thebibliography}

\end{CJK}
\end{document}
